% Options for packages loaded elsewhere
\PassOptionsToPackage{unicode}{hyperref}
\PassOptionsToPackage{hyphens}{url}
%
\documentclass[
]{article}
\usepackage{amsmath,amssymb}
\usepackage{lmodern}
\usepackage{graphicx}
\usepackage{iftex}
\ifPDFTeX
  \usepackage[T1]{fontenc}
  \usepackage[utf8]{inputenc}
  \usepackage{textcomp} % provide euro and other symbols
\else % if luatex or xetex
  \usepackage{unicode-math}
  \defaultfontfeatures{Scale=MatchLowercase}
  \defaultfontfeatures[\rmfamily]{Ligatures=TeX,Scale=1}
\fi
% Use upquote if available, for straight quotes in verbatim environments
\IfFileExists{upquote.sty}{\usepackage{upquote}}{}
\IfFileExists{microtype.sty}{% use microtype if available
  \usepackage[]{microtype}
  \UseMicrotypeSet[protrusion]{basicmath} % disable protrusion for tt fonts
}{}
\makeatletter
\@ifundefined{KOMAClassName}{% if non-KOMA class
  \IfFileExists{parskip.sty}{%
    \usepackage{parskip}
  }{% else
    \setlength{\parindent}{0pt}
    \setlength{\parskip}{6pt plus 2pt minus 1pt}}
}{% if KOMA class
  \KOMAoptions{parskip=half}}
\makeatother
\usepackage{xcolor}
\usepackage{longtable,booktabs,array}
\usepackage{calc} % for calculating minipage widths
% Correct order of tables after \paragraph or \subparagraph
\usepackage{etoolbox}
\makeatletter
\patchcmd\longtable{\par}{\if@noskipsec\mbox{}\fi\par}{}{}
\makeatother
% Allow footnotes in longtable head/foot
\IfFileExists{footnotehyper.sty}{\usepackage{footnotehyper}}{\usepackage{footnote}}
\makesavenoteenv{longtable}
\setlength{\emergencystretch}{3em} % prevent overfull lines
\providecommand{\tightlist}{%
  \setlength{\itemsep}{0pt}\setlength{\parskip}{0pt}}
\setcounter{secnumdepth}{-\maxdimen} % remove section numbering
\ifLuaTeX
  \usepackage{selnolig}  % disable illegal ligatures
\fi
\IfFileExists{bookmark.sty}{\usepackage{bookmark}}{\usepackage{hyperref}}
\IfFileExists{xurl.sty}{\usepackage{xurl}}{} % add URL line breaks if available
\urlstyle{same} % disable monospaced font for URLs
\hypersetup{
  hidelinks,
  pdfcreator={LaTeX via pandoc}}
\title{Replicación de índice con un portafolio de inversión utilizando clustering y programación Entera}
\author{Práxedis Jiménez Ruvalcaba}
\date{26 de mayo de 2025}

\begin{document}
\title{Replicación de índice con un portafolio de inversión utilizando clustering y programación Entera}
\author{Práxedis Jiménez Ruvalcaba}
\date{26 de mayo de 2025}
\textbf{Replicación de índice con un portafolio de inversión utilizando clustering y programación Entera}

Asesor: Beatriz Alejandra García Ramos

Práxedis Jiménez Ruvalcaba

\textbf{Introducción}

Lo que se busca lograr es un fondo indexado a través del
\emph{quantitative portfolio management} y, en particular, del uso de la
clusterización y la programación entera mixta para así conseguir
maximizar las ganancias y reducir los riesgos de este. En general se
busca adaptar el trabajo de Gharanchaei y Panda (2024)\footnote{Gharanchaei, M. K., Panda, P. P. (2024). \emph{Constructing and investment fund thourg stock clustering and integer programming}.Journal of Advancements in Applied Business Research.}, adaptado al
contexto particular de México.

No se busca un trabajo enfocado meramente a la inversión, sino que uno a través de la \emph{modern portfolio theory} ,en particular, el
trabajo realizado por Markowitz (1952)\footnote{Markowitz, H. (1952). \emph{Portfolio selection}. JSTOR. (pp 77- 91).}, y el marco en el que se constituye,
podamos realizar comparaciones y un análisis que nos permita un mayor
entendimiento de nuestra economía y la naturaleza de nuestro mercado
local. Lo cual, a su vez, nos puede facilitar comprender ciertos
eventos y condiciones históricas.

Es común hallar trabajos que toman datos de países alejados de la
realidad del nuestro, por lo que sus resultados también llegan a distar
de nuestro contexto. Esto nos incita a pensar que existe un valor en
adaptar otros trabajos para ver si son aplicables a nuestra realidad y
, si no, ampliar su obra en otros contextos.

\textbf{Revisión de literatura}

Con \emph{quantitative portfolio mangment}, nos referimos al uso de
herramientas de la matemática para optimizar la gestión de activos.
Siendo que este problema radica en dos; uno de clasificación (qué
activos considerar) y uno de optimización (cuánto de cada uno). Se
utilizarán dos herramientas, que son el \emph{clustering} y la programación
entera mixta.

Esta clase de herramientas ya han sido usado para este mismo tipo de
problemas, por ejemplo, en problemas como \emph{set covering} de Wolsey
(1998)\footnote{Wolsey Laurence, A. (1998). Formulations. \emph{Integer Programming}.(pp. 1-18). Wiley Interscience.}. Se utilizan conceptos de \emph{clustering} (al clasificar
conjuntos), en un problema de programación entera, solo que en este caso
siendo algo para minimizar costos.

De forma más específica, también ya se ha trabajado con optimizar
portafolios de inversión para máxima liquidez utilizando programación
entera mixta como en Bertsimas \emph{et al.} (1998)\footnote{Bertsimas, D., Darnell, C., Soucy, R., (1998). \emph{Portfolio construction through mixed integer programming}. MIT.}. Incluso ha habido
formas más recientes de trabajar con estas ideas utilizando herramientas
actuales como lo son optimizadores, por ejemplo, Gurobi Optimizer en
Kuhn (2023)\footnote{Kuhn, C. C. N., Calbert, G., Garanovich, I. L., Weir, T., (2023). \emph{Integer linear programming supporting portfolio design.} Elsevier}. Aunque en este caso era un problema con un enfoque
distinto al tener que ver más con administración gubernamental donde el enfásis radica en reducir riesgos, lo cual es muy distinto al usual querer maximizar ganacias.

Sin embargo, el trabajo que más servirá y más se cita como referencia
para esta clase de problemas en la literatura rescatada es sin duda
alguna, el trabajo de Markowitz (1952) de nombre \emph{mean-variance}
es hoy mejor conocido como \emph{modern portfolio theory}. Donde
describe y propone distintos modelos según ciertas propiedades, con el
objetivo de obtener los mayores retornos con los menores riesgos. En
particular describe tanto el \emph{equally weighted,} por capitalización
de mercado y de volatilidad mínima, además de la forma general de
optimizar.

Un fondo indexado, como lo indica su nombre, es un índice que monitorea
y mide los rendimientos de un conjunto particular de activos. Que en
este caso sería el índice de precios y cotización de la bolsa mexicana
de valores, que mide los rendimientos de 35 de las empresas más
importantes de México utilizando como métrica la capitalización de
mercado. La idea es tener un fondo con el que comprar activos imitando
los rendimientos de un índice que mide un mercado en particular. La cual
es una práctica común en casas de bolsa o fondos de inversión común.

La idea de realizar el trabajo con un fondo indexado en mente es que
estos son muy diversos y que se construyen en base a una estrategia
concreta, es decir, se elige un primer enfoque y el resto de las
decisiones son una consecuencia de esta primera estrategia seleccionada.
O, dicho de otra forma, es un enfoque pasivo. Lo cual dista de otros
tipos de portafolios donde se requiere estar tomando decisiones de forma
activa.

Existen trabajos que prueban una correlación entre maximizar las
ganancias y minimizar los riesgos como Anderson \emph{et al.} (2011)\footnote{Anderson, R. M., Bianachi, S. W., Goldberg, L. R. (2011). \emph{Will My Risk Parity Strategy Outperform?} Berkeley.} y Baker \emph{et al}. (2011)\footnote{Baker, M., Bradley, B., Wurgler, J\emph{. Benchmarks as Limits to Arbitrage: Understanding the Low-Volatility Anomaly}. Financial Analysts Journal.}, de forma que al realizar uno de los
objetivos, implícitamente estás maximizando riesgos o minimizando
ganancias. Dichos trabajos remarcan que existe evidencia para creer que
mientras se minimicen los riegos un efecto es la reducción de las
ganancias.

Como ninguno de los enfoques sobre si maximizar ganancias afecta a los
riesgos parece tener evidencia para descartar al otro. Además de que en
los resultados de estos trabajos y nuestras referencias principales;
Gharanchaei y Panda (2024), Markowitz (1952), muestran que maximizar
ganancias minimiza el riesgo en cierta medida, pero esto se puede
optimizar. Tendremos esta visión a la hora de trabajar con el problema.

El tomar la correlación como métrica para la clusterización tiene un
fundamento teórico en economía, pues acorde a Koumou (2020)\footnote{Koumou, G. B.(2020) \emph{Diversification and portfolio theory: a review}. Swiss Society for Financial Market Research.}, una de las
maneras de diseñar la correcta diversificación de un portafolio es
agrupando a los activos con una correlación cercana a 1 y separando al
resto. Esto reduce las probabilidades de que distintas agrupaciones
tengan un mal desempeño al mismo tiempo, a esto se le conoce como
\emph{multilateral insurance}.

Tomamos la diversificación como algo bueno ya que, si bien Choueifaty y
Coignard (2008) \footnote{Choueifaty, Y., Coignard, Y. (2008). \emph{Toward Maximum Diversification}. The Journal of Portfolio Management.} muestran que existen diversos enfoques donde puede ser
una estrategia no óptima, también se dan argumentos contrarios y que es
más prominente las posibles ventajas que las desventajas. Además de que
se remarca que los modelos no necesariamente reflejan la realidad del
mercado, aunque puedan llegar a interpretarlo apropiadamente, y el éxito
que pueda tener un modelo no debe de ser el único criterio para
trasponerlo en el mundo real.

En nuestro caso, al estar imitando un índice, lo que buscamos es
precisamente un modelo que refleja la realidad de un mercado y economía
específicos. Esto por el interés de que todas las consideraciones en los
trabajos citados son en un mercado competitivo y diverso, siendo esto
algo distinto a la realidad del mercado mexicano el cual es monopólico y
poco diverso. (Reed \emph{et al,} 2023)\footnote{Reed, T; Muñoz Moreno, R. (febrero 23, 2023). \emph{How the Federal Economic Competition Commission fights for workers and consumers in Mexico}. World Bank.}.

Choueifaty y Coignard (2008) enfatizan que no hay una sola medición que
sirva para determinar la optimalidad de un modelo y es por ello por lo
que todo trabajo de índole bursátil, debe de contar con al menos dos
enfoques distintos. Por ejemplo, existen argumentos válidos que
favorecen a la varianza entre activos como una medida, pero también los
hay que dicen que el interés compuesto es una mejor forma de determinar
qué tan bien funciona el modelo.

En nuestro caso estaremos tomando el enfoque de las estrategias
\emph{equally weighted} ya que por Chow \emph{et al.} (2011)\footnote{ Chow, T., Hsu, J., Kalesnik, V., Little, B., \emph{A Survey of Alternative Equity Index Strategies}. Financial Analysts Journal.}, y Carvalho \emph{et al.} (2012) \footnote{  Carvalho R. L., Lu, X., Moulin, P. (2012).\emph{Demystifying Equity Risk-Based Strategies: A Simple Alpha plus Beta Description}. The Journal of Portfolio Management.}, entendemos que, si bien se trata de una estrategia
que antepone minimizar el riesgo a las ganancias, sigue siendo óptima
además de que puede ser adaptada para asemejarse a otros tipos de
estrategias que buscan diversificar, las cuales pueden llegar a ser más
eficiente pero solo cuando se trata de reducir los costos, especialmente
durante la inversión inicial.

A estas estrategias se les llama también \emph{risk parity assest
allocation} (Choueifaty y Coignard, 2008), que nuevamente es distribuir
los ``pesos'' para repartir el riesgo de manera equitativa. Siendo que
el riesgo tiende a estar medido en volatilidad. Donde la volatilidad
puede entenderse como qué tanto sube y baja de precio, por lo que
podemos medirla como desviación estándar de los precios de una acción en
un cierto período de tiempo.

Baker \emph{et al.} (2011), muestra que existen otras formas de medir la
volatilidad dependiendo del tipo de activo y las consideraciones que se
busquen tener, pero no es que exista una mejor forma que otra, sino que
depende del contexto, aunque las formas más básicas tienden siempre a
dar un acercamiento bastante acertado y por ello tomaremos la
volatilidad como la desviación estándar de los precios de cierre en el
período del tiempo de reajuste de los portafolios.

En ese mismo trabajo también se concluye que los portafolios
diversificados son óptimos en comparación del resto, aun considerando un
horizonte de tiempo real y distintas métricas. Koumou. (2020), Da
argumentos a favor y en contra de la diversificación, pero su conclusión
es también que la diversificación es una estrategia óptima. Por ello
este es el enfoque en el que estamos abordando el problema.

Otra métrica que utilizaremos es la de la capitalización de mercado, la
cual son las acciones en circulación de una empresa particular
multiplicada por el valor que tienen dichas acciones de forma individual
en dicho momento. Entonces la capitalización de mercado no es más que el
valor total de las acciones de circulación.

La capitalización de mercado es una métrica muy usada porque te habla
sobre la presencia que una empresa tiene en un mercado, entre más alta
sea, más presencia tiene. A su vez, las fluctuaciones que pueda tener en
un cierto tiempo te hablan también sobre la volatilidad que dicha
empresa. Ambas lecturas dan información tanto de si una empresa es
representativa de un sector como del riesgo que pueda significar
invertir en esta, por lo tanto, es de interés para nosotros además de
que el índice que estamos intentando replicar esta medido con esta
métrica.

El algoritmo que se quiere usar para el clustering es el de
\emph{partitioning around medoids}, Botyarov y Miller (2022)\footnote{Botyarov, M., Miller, E. E., \emph{Partitioning around medoids as a systemic approach to generative design solution space reduction.} Elsevier.}; el cuál es
semejante a \emph{K-means} por tener una cantidad fija de centroides,
pero sin usar la media aritmética y distancias euclidianas para
relacionar a los elementos a un centroide, sino que, usando cualquiera
otra métrica para medir la semejanza de los elementos con el centroide,
que en este caso será la correlación.

En algunos trabajos los activos ya son un conjunto de otros activos que
poseen una alta similitud entre sí, esto es solo para hacer más sencillo
el seguimiento del desempeño de los activos, por el trabajo que se debe
hacer de estar constantemente monitoreando y analizando los desempeños
de estos activos, por lo que cada activo asignado a un clúster es
representativo de otro. Bajo esta misma lógica es que la cantidad de
centroides se elige notablemente menor a la de activos totales
(Gharanchaei y Panda, 2024).

\textbf{Descripción del problema.}

Hay dos enfoques con los que abordar el problema de replicar un índice;
heurísticas y optimización. Las heurísticas es con las 3 fórmulas que
serán mencionados a continuación del modelo para el clustering, pero se
consiguen con base a la teoría y utilizan el concepto de
\emph{multilateral insurance,} así como las métricas de la
capitalización de mercado y la volatilidad.

Mientras que el enfoque por optimización consiste en utilizar el trabajo
de Markowitz (1952), y usar sus algoritmos para mínima varianza y
máxima ganancia, pero en una forma matricial y aprovechando los ya
resultados proporcionados por las heurísticas, que en este caso nos
estaremos refiriendo como \emph{benchmarks}, las cuales serán las
optimizadas.

El índice que intentaremos replicar con lo antes mencionado y
posteriormente descrito es el S\&P índice de precios y cotizaciones de
la bolsa mexicana de valores (S\&P IPC BMV), comúnmente abreviado como
IPC.

Tomamos este índice ya que es uno de los principales indicadores de la
economía mexicana, pues su comportamiento ha descrito apropiadamente el
de nuestra economía y es una buena muestra de los principales sectores
económicos de nuestro país. De mayor a menor proporción en el IPC se
encuentran los sectores de servicio de telecomunicaciones, consumo
frecuente, materiales, servicios financieros, industrial, servicio y
bienes de consumo no básicos y salud.

El IPC es elaborado por la compañía S\&P Dow Jones, la cual está
enfocada en realizar índices de mercado. En este caso el IPC se elabora
utilizando como métrica la capitalización de mercado de 35 empresas
mexicanas S\&P Down Jones (2025)\footnote{S\&P Dow Jones Indices. (31 marzo 2025). S\&P BMV IPC: \emph{Constituents Export}. {[}Base de datos{]}.}, las cuales
son:

\begin{longtable}[]{@{}
  >{\raggedright\arraybackslash}p{(\columnwidth - 2\tabcolsep) * \real{0.5000}}
  >{\raggedright\arraybackslash}p{(\columnwidth - 2\tabcolsep) * \real{0.5000}}@{}}
\toprule()
\multicolumn{2}{@{}>{\raggedright\arraybackslash}p{(\columnwidth - 2\tabcolsep) * \real{1.0000} + 2\tabcolsep}@{}}{%
\begin{minipage}[b]{\linewidth}\raggedright
\textbf{Empresas mexicanas parte del IPC}
\end{minipage}} \\
\midrule()
\endhead
Alfa S.A. A & Grupo Bimbo S.A.B. \\
Alsea S.A. & Grupo Carso S.A.B. de C.V. \\
América Móvil S.A.B. de C.V. B & Grupo Cementos de Chihuahua S.A.B. de
C.V. \\
Arca Continental S.A.B. de C.V. & Grupo Comercial Chedraui S.A. de
C.V. \\
Banco del Bajío S.A. & Grupo Financiero Banorte O. \\
Becle S.A. De C.V. & Grupo Financiero Inbursa O \\
Bolsa Mexicana de Valores S.A. de C.V. & Grupo México S.A.B. de C.V.
B \\
Cemex S.A. CPO & Grupo Televisa S.A.B. CPO \\
Coca Cola Femsa S.A.B. de C.V. UBL & Industrias Peñoles \\
Corporación Inmobiliaria Vesta S.A.B. de C.V. & Kimberly Clark de México
S.A.B. de C.V. A \\
El Puerto de Liverpool S.A.B. de C.V. & La Comer S.A.B. de C.V. UBC \\
Fomento Económico Mexicano S.A.B. de C.V. & Megacable Holdings S.A.B. de
C.V. \\
Genomma Lab Internacional S.A. de C.V. & Orbia Advance Corporation
S.A.B. de C.V. \\
Gentera S.A.B. de C.V. & Promotora y Operadora de Infraestructura SAB de
CV \\
Gruma S.A.B. B & Qualitas Controladora S.A.B. de C.V. \\
Grupo Aeroportuario del Centro Norte S.A.B. de C.V. & Regional S.A. de
C.V. \\
Grupo Aeroportuario del Pacifico, S.A.B. de C.V. & Walmart de México
S.A.B. de C.V. \\
Grupo Aeroportuario del Sureste S.A.B. de C.V. B & \\
\bottomrule()
\caption{Listado de empresas que conforman el IPC\&P BMV, acorde a la información obtenida proporcionada por los elaboradores del ínidce S\&P Down Jones.}
\end{longtable}

Utilizando las librerías ``yfinance'' e ``invest-corepy'' en Visual
Basic 1.99.3 con Python 3.12.7 integrado en el paquete de anaconda de
versión 2.6.5, se obtuvieron los datos históricos desde Yahoo Finance de los precios máximos, mínimos y
de cierre, además de la cantidad en circulación de las acciones de las
35 empresas durante los períodos 2008 a 2010 y 2015 a 2024.

Se tomaron estas fechas ya que se desea probar cómo funcionan los
distintos enfoques previamente mencionados en un período de diez años,
específicamente el de 2015 a 2024 ya que en 2010 se cambiaron algunas
consideraciones al momento de medir el IPC, así que esas fechas poseen
una metodología consistente de mediciones en el índice, además de que
son fechas con crecimiento económico moderado, una recesión en 2020 y
una recuperación en 2021 (Banco de México, 2025) \footnote{Banco de México. (2025). Sistema de Información Económica {[}Base de datos{]}.}. 

Podemos decir que estamos tomando un período nada despreciable para
estudiar, dado que todo tipo de evento en la economía tomó lugar en este
período y en particular para realizar un contraste con Garchaei y Panda
(2024), donde se toma el período 2014 a 2023.

Además, sería ideal probar los modelos en el período de 2008 a 2010,
para ver cómo contrasta con una recesión mucho mayor a la del 2020 y
verificar qué tan robusto es el modelo. En la literatura se ha visto una
práctica común y útil el probar modelos durante las crisis de 1929 y
2008, puesto que han sido de las más graves para Estados Unidos y ser
capaz de tener poco error entonces indica una buena robustez para
modelos o una buena certidumbre conforme la teoría que se esté probando.

La fecha particular de 2008 no es solo porque se haga en trabajos
enfocados a Estados Unidos, sino que, en el caso de la economía
mexicana, hemos padecido de 5 recesiones importantes posteriores a la
guerra cristera de 1926; la de 1980, 1994, 2000, 2008 y 2020. Pero
tomamos 2008 porque es donde las condiciones de mercado son más
similares a las actuales y no hay que ampliar mucho en el marco del contexto
histórico a la hora de realizar un análisis.

Cabe destacar que para 2008, los datos de Banco del Bajío S.A., Becle
S.A. de C.V., Coca Cola Femsa S.A.B. de C.V. UBL, La Comer S.A.B. de
C.V. UBC, Qualitas Controladora S.A.B. de C.V. y Regional S.A. de C.V.,
están incompletos y no fue posible completarlos.

Los portafolios, excepto el del clustering y el \emph{equally weighted}
que por su naturaleza son fijos, van a estarse reajustando cada 6 meses.
Lo ideal es tener distintos períodos para reajustarse, pero tomamos los
6 meses porque es la misma cantidad de tiempo en la que IPC se
requilibra, entonces así tenemos una cierta consistencia entre nuestro
portafolio y aquello que queremos imitar.

Los reajustes semestrales para la primera mitad serán partiendo del 30 de
junio para los años 2015, 2016, 2017, 2020, 2021, 2022 y 2023. Del 28 de
junio de para 2019 y 2024. Y del 29 de junio para 2018. Mientras que para la
segunda son el 31 para los años 2015, 2018, 2019, 2020 y 2021. 30 de
diciembre para 2016, 2022, 2024. Y 29 de diciembre para 2017 y 2023.

La primera parte del problema es obtener las correlaciones lineales
entre acciones, puesto que será la métrica que usaremos para nuestra
clusterización usando \emph{partiotoning around medoids}. Esto nos
indicará cuáles son las acciones más representativas de cada clúster.

Seleccionamos la cantidad de 2, 4 y 8 clústers para realizar el
\emph{partitioning around medoids}, con base a proporciones acorde al
trabajo de Gharanchaei y Panda (2024), ya que su trabajo considera 84
series de acciones y 5, 10 y 20 clústeres. Como nosotros consideramos
solo 35 series de acciones, entonces estimamos las proporciones y así
establecimos las cantidades de clústeres con las cuales trabajar el
problema.

La manera en que queremos optimizar la culsterización es con uso de
programación entera mixta, pues (Gharanchaei y Panda., 2024)
describe un modelo con el cual contamos con variables enteras, pero
nuestra función objetivo maximiza un número real el cual es la suma de
las correlaciones de los elementos de los activos con su respectivo
centroide.

\begin{equation}
\max\sum_{i = 1}^{n}{\sum_{j = 1}^{n}{c_{ij}x_{ij}}} \tag{1}
\end{equation}
Sujeto a:

\begin{equation}
\sum_{j = 1}^{n}{y_{j} = k} \tag{1.2}
\end{equation}
\begin{equation}
\sum_{j = 1}^{n}{x_{ij} = 1\ \forall i} \tag{1.3} 
\end{equation}
\begin{equation}
x_{ij} \leq y_{j}\ \forall i,j \tag{1.4}
\end{equation}
\begin{equation}
x_{ij},y_{j}\epsilon\left\{ 0,1 \right\} \tag{1.5}
\end{equation}


Donde las variables toman sus valores acordes a lo descrito en la
siguiente tabla:

\begin{longtable}[]{@{}
  >{\raggedright\arraybackslash}p{(\columnwidth - 2\tabcolsep) * \real{0.2195}}
  >{\raggedright\arraybackslash}p{(\columnwidth - 2\tabcolsep) * \real{0.7805}}@{}}
\toprule()
\begin{minipage}[b]{\linewidth}\raggedright
\[c_{ij}\]
\end{minipage} & \begin{minipage}[b]{\linewidth}\raggedright
Similitud entre activo i y j. Es la correlación que existe entre ambos
activos. Es un valor entre -1 y 1.
\end{minipage} \\
\midrule()
\endhead
\(x_{ij}\) & Indica si el activo j del portafolio posee una alta
correlación positiva al activo i. Toma el valor de 1 si este es el caso
y toma 0 de lo contrario. \\
\midrule()
\(y_{j}\) & Denota qué activo j se encuentra en el portafolio tomando el
valor de 1, de lo contrario es 0. \\
\midrule()
\emph{k} & La cantidad de centroides, el cual como ya se explicó, debe
de ser un número más chico que el total de activos porque se trabajará
con representantes de la clase. \\
\bottomrule()
\caption{Descripción de las variables utilizadas para realizar el \textit{clustering} de \textit{k-medoids} para obtener activos representativos.}
\end{longtable}

Ya tomando solo las acciones representativas, tenemos con qué medir el
error, puesto que es como tomar los principales componentes el IPC. Y
ahora podemos trabajar con el enfoque para resolver el problema a través
de \emph{benchmarks}, ya anteriormente mencionados en la revisión de
literatura y que tomamos de Carvalho \emph{at al.} (2012). Aunque son
también mencionados por Markowitz (1952).

La ecuación (1) maximiza el producto de la correlación lineal, para así poder seleccionar aquellos que posean mayor correlación con el resto y que sean una muestra representativa. La restricción (1.2) nos garantiza que se asignen los mismos activos que centroides. Mientras que las restricciones (1.3) y (1.4) son para evitar que se puedan repetir la selección de ciertos activos, que al trabajar con una matriz cuadrada puede que se seleccionen el mismo \emph{x} \emph{i,j} y el \emph{j,i} al mismo tiempo, lo cual no sería una solución real. (1.5) solo establece la naturaleza de las variables como binarias.

Para mantener una consistencia con entre los \emph{benchmarks} y el
enfoque por optimización del trabajo de Markowitz (1952), donde se
destaca el considerar el problema sin tomar en cuenta la participación o
las posibles decisiones que un asesor, corredor de bolsa o la
experiencia y criterio que alguien pueda aportar, especialmente en
cuestión probabilística, lo descartaremos. Solo tomaremos en cuenta los
resultados directos.

Lo primero que haremos será un portafolio de tipo \emph{equally
weighted.} El cual es simplemente darle el mismo peso a todas las
acciones para poder tener resultados consistentes con lo anteriormente
planteado y poder contrastar resultados, tomaremos 3 de estos modelos,
empezando por el más general y mencionado:

Si tienes \emph{n} activos
\begin{equation}
 w_{i} = \frac{1}{n} \tag{2}   
\end{equation}

\begin{longtable}[]{@{}
  >{\raggedright\arraybackslash}p{(\columnwidth - 2\tabcolsep) * \real{0.2082}}
  >{\raggedright\arraybackslash}p{(\columnwidth - 2\tabcolsep) * \real{0.7918}}@{}}
\toprule()
\begin{minipage}[b]{\linewidth}\raggedright
\[w_{i}\]
\end{minipage} & \begin{minipage}[b]{\linewidth}\raggedright
Peso del activo \emph{i}, es decir, su importancia y cuánto debe de
invertírsele en comparación del resto
\end{minipage} \\
\midrule()
\endhead
\(n\) & Total de activos \\
\bottomrule()
\caption{Descricpión de variables para realizar el \textit{benchmark equally weighted.}}
\end{longtable}

Este enfoque consigue una diversificación tal que maximiza la
rentabilidad de la inversión con relación al riesgo que poseen los
activos que se asumen. Simplemente es otorgarles el mismo peso o
prioridad a todos los activos, que por la anteriormente
\emph{multilateral insurance}, termina por proporcionar un riesgo mínimo
sin condicionar la ganancia máxima.

\emph{Equally Wighted} es el \emph{benchmark} más noble. Es simple de
describir, operar, no requiere estarlo requilibrando en ningún momento.
Y el no desplomar la ganancia a costa de reducir el riesgo, además de
servir en cualquier mercado y ser fácilmente adaptable; lo ha vuelto una
estrategia vencedora difícil de superar.

El otro modelo sería el del \emph{market-cap}, el cual asocia cada
activo acorde a su capitalización del mercado, el cual es el producto de
acciones en circulación y su precio actual.

\begin{equation}
w_{i} = \frac{M_{i}}{\sum_{j}^{n}M_{j}}   \tag{3}
\end{equation}

Donde

\begin{equation}
 M_{i} = x_{i}p_{i} \tag{4}   
\end{equation}

\begin{longtable}[]{@{}
  >{\raggedright\arraybackslash}p{(\columnwidth - 2\tabcolsep) * \real{0.1601}}
  >{\raggedright\arraybackslash}p{(\columnwidth - 2\tabcolsep) * \real{0.8399}}@{}}
\toprule()
\begin{minipage}[b]{\linewidth}\raggedright
\[M_{i}\]
\end{minipage} & \begin{minipage}[b]{\linewidth}\raggedright
Capitalización del mercado de la compañía \emph{i.} Puede entenderse
como una métrica con la cual evaluar la viabilidad de invertir en una
compañía. Cocientando esta medida con la suma de todas podríamos saber
cuál es la proporción con la que debemos invertir en cada compañía.
\end{minipage} \\
\midrule()
\endhead
\(x_{i}\) & Cantidad de acciones circulando en el mercado de la compañía
\emph{i} \\
\midrule()
\(p_{i}\) & Precio actual de las acciones de la compañía \emph{i} \\
\bottomrule()
\caption{Descricpión de variables del \textit{benchmark} de la capitalización de mercado.}
\end{longtable}

Esta es la estrategia de este tipo de enfoque que menos genera. Es en
cierta forma muy segura y cambia muy poco, pues al generar poco, uno
tarda en reinvertir los dividendos generados por la inversión anterior.

Lo que plantea el \emph{benchmark} es invertir proporcionalmente a la
capitalización de mercado, entre más alta sea, más habrá que invertir.
Ya que esto asegura estar invirtiendo acorde a las propias proporciones
del mercado y así tomar solo acciones representativas por lo que reduces
el riesgo. Tener activos en las empresas más grandes de cada sector es
considerado como una estrategia a veces excesivamente segura, de ahí
viene su usual poca ganancia, pero gran reducción de riesgo.

Al momento de medir la capitalización de mercado los datos que tenemos
son las acciones en circulación en un cierto momento. No disponemos de
la información necesaria para saber si algunas de esas acciones son
\emph{inside trading,} es decir, al intercambio privado entre
accionistas o miembros de la compañía de acciones\emph{.} Por lo que
consideraremos que todo lo que está en circulación es abierto para todo
el público o todo el mercado accionario al menos.

La ecuación (3) parte de distribuir los pesos directamente proporcional a la capitalización de mercado. Entre mayor sea su capitalización de mercado, mayor peso. Esto ya que entre más activos en circulación haya de un activo, hay menos probabilidades de que lo vendan o lo compren todo, además de que como es multiplicado por el valor del activo, te aseguras de que las acciones posean una cierta demanda y no vayan a perder su precio fácil al haber muchos en circulación. El único problema que tiene es si el volumen es tan grande o si el precio es tan grande que el resto de componentes dejen de importar, pero esto es inusual en una economía competitiva.

El último enfoque del tipo de diversificación que estamos trabajando y
usaremos es el que se basa en la volatilidad como métrica. Reiterando
algo mencionado previamente, en este caso mediremos la volatilidad con
la desviación estándar de los precios de un activo. El modelo en
cuestión es:

\begin{equation}
w_{i} = \frac{\frac{1}{\sigma_{i}}}{\sum_{j}^{n}\frac{1}{\sigma_{j}}}   \tag{5}  
\end{equation}

\begin{longtable}[]{@{}
  >{\raggedright\arraybackslash}p{(\columnwidth - 2\tabcolsep) * \real{0.1440}}
  >{\raggedright\arraybackslash}p{(\columnwidth - 2\tabcolsep) * \real{0.8560}}@{}}
\toprule()
\begin{minipage}[b]{\linewidth}\raggedright
\[\sigma_{i}\]
\end{minipage} & \begin{minipage}[b]{\linewidth}\raggedright
Volatilidad del activo \emph{i}. La desviación estándar de su precio.
\end{minipage} \\
\endhead
\bottomrule()
\caption{Descripción de las variables del \textit{benchmark} de la volatilidad.}
\end{longtable}

Se trata de un caso particular puesto que, en condiciones concretas,
llega a ser de lejos el más eficiente de los tres modelos, pero en general
se trata de un enfoque centrado en favorecer los activos de bajo riesgo,
pero teniendo rendimientos más similares al primer modelo y de hecho
siendo a veces muy similar a este.

La manera en que funciona es simplemente con proporciones, como queda indicado en (5). El activo con
mayor riesgo es el que menos peso deba tener en el fondo indexado,
mientras que el de menor riesgo va a tener el mayor peso. Esto nos da
una certidumbre al momento de minimizar el riesgo, pero lo consigue
llegando a reducir considerablemente las ganancias.

Para asignar el peso acorde a la capitalización de mercado de los
clústeres, tenemos la fórmula:

\begin{equation}
 x_{jl} = \frac{\sum_{i \in C_{l}}^{}V_{i}}{\sum_{I = 1}^{n}V_{i}},\ \ l = 1,\ldots,\ k  \tag{6} 
\end{equation}

Donde \emph{k} es el total de los clústeres

\begin{longtable}[]{@{}
  >{\raggedright\arraybackslash}p{(\columnwidth - 2\tabcolsep) * \real{0.5000}}
  >{\raggedright\arraybackslash}p{(\columnwidth - 2\tabcolsep) * \real{0.5000}}@{}}
\toprule()
\begin{minipage}[b]{\linewidth}\raggedright
\[V_{i}\]
\end{minipage} & \begin{minipage}[b]{\linewidth}\raggedright
Capitalización de mercado de la acción i
\end{minipage} \\
\midrule()
\endhead
\(C_{i}\) & Clúster i \\
\midrule()
\emph{l} & Cantidad de clústeres \\
\midrule()
\(x_{jl}\) & Peso de la acción j en el clúster l \\
\bottomrule()
\caption{Descripción de variables para asignar pesos optimizados según los activos representativos de la clusterización y la capitalización de mercado.}
\end{longtable}

Lo que se realiza aquí es determinar la proporción que cada acción
debería de tener en cada uno de los clústeres que se determinaron al
inicio, acorde a la capitalización de mercado. Por eso es la
capitalización de mercado de un cierto clúster, entre la del total de
clústeres.

 Las ecuaciones (2), (3), (4) y (5), no se obtienen de manera experimental, sino que se determinan por planteamiento teórico. A través de la propia teoría de la economía se determinan estas ecuaciones. Uno puede fijarse en que simplemente son determinar una proporción acorde a una cierte métrica del rendimiento de un activo, en estos casos con el objetivo de minimizar el riesgo. Ya sea repartiendo con la seguridad que da el mercado competitivo con la capitalización de mercado, la seguridad que da darle menos peso a lo que tiene más riesgo con la volatilidad, o la \emph{multilateral isnurance} con la repartición uniforme de pesos.

Cabe resaltar que el uso de la capitalización de mercado como factor de
decisión para determinar las proporciones es porque es la métrica con la
que se elabora el índice que deseamos imitar. El IPC es elaborado con la
capitalización de mercado por lo que usar estás métricas durante el
proceso de optimización es en parte para mantener una consistencia y
esperando conseguir un mejor resultado, además de que nos permitirá
realizar ciertos análisis de forma sencilla, sin tener que estar tomando
a consideración equivalencias y nuevamente, aprovechar una cierta
consistencia.

Este proceso de clusterización debe de realizarse para poder trabajar
apropiadamente con el siguiente modelo de optimización, el cuál debería
de optimizar la varianza en base a la ganancia máxima obtenida en los
anteriores modelos. Se usará el siguiente modelo para cada uno de los tres
\emph{benchmarks} anteriormente planteados.

Estos serían los tres modelos que tomaríamos como referencia para
nuestros índices y ver cómo se comporta la clusterización realizada en
base a la correlación. Para optimizar los portafolios anteriormente
realizados con los \emph{benchmarks} tomaríamos como referencia el
siguiente modelo planteado por Choueifaty y Coignard (2008):

\begin{equation}
\min{\left( X - X_{B} \right)^{T}M(X - X_{B}}) \tag{7}  
\end{equation}

Sujeto a:

\begin{equation}
1^{T}X = 1 \tag{7.2}    
\end{equation}
\begin{equation}
X \geq 0 \tag{7.3}    
\end{equation}
\begin{equation}
\left| x_{t} - x_{t - 1} \right| < cota\ de\ ganancia\ máxima \tag{7.4}
\end{equation}
\begin{equation}
\sum_{i = 1}^{n}{\left| x_{i - 1} - x_{i} \right| \leq t\ para\ t\ que\ es\ un\ "treshold"}   \tag{7.5}
\end{equation}
\begin{equation}
x_{i}\epsilon\left\{ 0,1 \right\}\ para\ i = 1,\ldots,n    
\end{equation}

\begin{longtable}[]{@{}
  >{\raggedright\arraybackslash}p{(\columnwidth - 2\tabcolsep) * \real{0.1601}}
  >{\raggedright\arraybackslash}p{(\columnwidth - 2\tabcolsep) * \real{0.8399}}@{}}
\toprule()
\begin{minipage}[b]{\linewidth}\raggedright
\emph{X}
\end{minipage} & \begin{minipage}[b]{\linewidth}\raggedright
Es un vector con las acciones del portafolio creado con clustering donde
\(x_{i}\) es el activo \emph{i.}
\end{minipage} \\
\midrule()
\endhead
\(X_{b}\) & Es un vector con los pesos asignados a un portafolio de los
creados con los modelos donde \(x_{b}\) es el activo \emph{b} que
corresponde con el \emph{i}, es decir, son el mismo activo de la misma
compañía. \\
\midrule()
\emph{M} & Es el vector con la capitalización de mercado de cada activo
\emph{i.} \\
\bottomrule()
\caption{Descripción de variables para minimización de error utilizando teoría de Markowitz.}
\end{longtable}

La idea es que la métrica sea el error y las ganancias. En general la
manera en que se miden esto y que se consigue este modelo es a través de
heurísticas, pero el error sigue funcionando de forma usual, es decir,
midiendo la desviación estándar de la diferencia de ganancia de los
modelos y lo que se tiene en el período previo. Se intenta balancear
esta varianza con las diferencias en los coeficientes de los portafolios
de los modelos y el del clustering.

La ecuación (7) es una que planteó Markowitz (1952) para hallar el riesgo mínimo dado unos ciertos rangos de ganancia. En este caso la ecuación es planteada con matrices dado que es una forma de plantear sistemas de ecuaciones con varias variables. Los rangos de ganacias estarían dados por la diferencia entre lo obtenido por los \emph{benchmark} y el valor del conjunto l de los activos que llevamos en el portafolio (aquello que queremos determinar). En cierta medida es como minizar el error para obtener un valor esperado. 

La restricción (7.2) nos garantiza que se asignen apropiadamente los valores y no nos quedé un vector vacío o que el modelo de resultados falsos o no viables. La restricción (7.3) se base en que no se considerará las ventas inmediatas de activos. Durante el proceso se tendrá o no un cierto activo, no se tomará en cuenta el caso de venderlo.
La restricción (7.4) es para tener una cota de ganancia máxima y estar obteniendo el riesgo mínimo acorde a esa cota. Esta restricción es necesaria cuando se trabaja en más de dos dimensiones puesto que el espacio de soluciones puede no ser acotado y entonces ningún modelo daría una solución óptima. La última restricción (7.5) es para tomar en cuenta los períodos de rebalanceo del portafolio y no realizar cambios que puedan presentar un problema a largo plazo.

Una vez los resultados de los modelos sean obtenidos, entonces mediremos
el error para analizar si en verdad la teoría y esta propuesta es apta
para una economía como la mexicana. Para los \emph{benchmarks} la
mediremos obteniendo la desviación estándar de la diferencia entre lo
que el \emph{benchmark} y el portafolio de la clusterización hayan
ganado.

\textbf{Experimentación y análisis de resultados}

Todo fue realizado en una laptop con sistema operativo Microsoft Windows
11 Home Single Language, Versión 10.0.26100, de marca Acer y modelo
Aspire AV15-53P, con un procesador de décima tercera generación Intel
Core(TM) i5-1335U,1300 Mhz. Una memoria física instalada (RAM) de 16 GB,
y un BIOS Insyde Corp V1.09.

\textbf{Generación de instancias}

El código fue realizado en un jupyter notebook de Python, de manera que cada casilla correspondía a una parte distinta del proceso. Primero se tuvo que elaborar el código que utilizando las librerías de “yfinance” e “invest-corepy” para poder extraer los datos diarios de las acciones de cada una de las 35 empresas del IPC listadas anteriormente. 

La información de cada empresa era organizada en un dataframe con las columnas que indicaban la fecha, el precio de cierre de la acción, el precio máximo, el mínimo y la cantidad de acciones en circulación. Se ordenaban todas en ese preciso orden para posteriormente guardarse en un archivo CSV. 

Después se volvían a leer los archivos y se iban guardando los datos en un mismo \emph{dataframe}, solo que organizando los nombres de cada columna de manera que indicaban qué dato eran y de cuál empresa. Cabe resaltar que todos los datos tenían las mismas fechas, por la que no había problemas al momento de ir concatenando los \emph{dataframes}. Además de esto se guardaban dos copias del \emph{dataframe}, excepto que una solo poseía las columnas que indicaban el precio de cierre.

Es con este último es que se obtenían las correlaciones lineales entre los precios de todos y con dichos datos hacíamos el clustering haciendo uso del modelo (1), aprovechando el optimizador Gurobi (versión 12.0.0). Esto nos daba un archivo de texto el cual indicaba los activos representativos por clúster y a qué clúster eran asignados cada otro activo. Los activos representativos los escribía manualmente en una lista que se guardaría en memoria. También de forma manual se hacía una lista que guardaba la separación de cada semestre, el día en el que iniciaba uno y cuándo terminaba otro. Además de otra que hacía lo mismo para los datos del IPC, los cuales eran mensuales.

También se indexaba la copia del \emph{dataframe} que no se había utilizado, con las fechas como índice, esto para la siguiente celda que se encargaba de conseguir los pesos que debía tener cada uno de los activos seleccionado a través del \emph{clustering}, durante cada semestre. De manera que se guardaba todo en una lista de listas, que poseía el peso de cada activo durante un semestre determinado. Se guardaba una lista de los pesos acorde a la capitalización de mercado, utilizando la ecuación (3) y otra con los de la volatilidad usando la (5).

Cabe resaltar que se debía señalar de forma manual la cantidad de clústeres, para que el programa pudiese determinase los pesos. No era necesario especificar las empresas, pero sí la cantidad de clústeres. Ya después de esto se calculaba el error en otra celda. Pero primero hay que destacar que para el \emph{benchmark equally weighted} (2), no era necesario realizar un cálculo más que obtener el inverso multiplicativo del total de activos, que en este caso se trataba de la cantidad de clústeres. 

La última celda era calcular el error con la desviación estándar de la diferencia del valor del portafolio al término del semestre con los pesos que determinaban los \emph{benchmark}, con el valor del IPC. También se midió el error tomando cambiando el valor explicito del IPC, por el de la suma de las capitalizaciones de mercado de todas las empresas que conforman el índice, es decir, una medición proporcional del IPC. Al final se consideró esta última ya que no se había considerado el cambio de empresas en la composición y métricas del índice, por lo que esta medida era mucho más congruente con el trabajo que habíamos realizado.

Además del cálculo, los resultados eran guardados en un CSV que después manualmente se convirtió en un archivo de Excel, para poder analizar los resultados detenidamente y poder realizar gráficas de manera más cómoda. 

\textbf{Resultados}

Las siguientes tablas representan los resultados obtenidos por el clustering, cada tabla es el resultado de una selección distinta de centroides. Recordemos que decidimos tomar 2,4 y 8 centroides.

\begin{longtable}[]{@{}
  >{\raggedright\arraybackslash}p{(\columnwidth - 2\tabcolsep) * \real{0.4973}}
  >{\raggedright\arraybackslash}p{(\columnwidth - 2\tabcolsep) * \real{0.5027}}@{}}
\toprule()
\multicolumn{2}{@{}>{\raggedright\arraybackslash}p{(\columnwidth - 2\tabcolsep) * \real{1.0000} + 2\tabcolsep}@{}}{%
\begin{minipage}[b]{\linewidth}\raggedright
\textbf{Empresas representativas para 2 clústeres}
\end{minipage}} \\
\midrule()
\endhead
\textbf{Grupo Financiero Banorte O.} & \textbf{Grupo Televisa S.A.B. CPO} \\
America Movil S.A.B. de C.V. B & Alfa S.A. A. \\
Arca Continental S.A.B. de C.V. &  Becle S.A. De C.V. \\
Grupo Bimbo S.A.B. & Bolsa Mexicana de Valores S.A. de C.V. \\
Cemex S.A. CPO & Industrias Peñoles \\
Banco del Bajio S.A. & Megacable Holdings S.A.B. de C.V.\\
Alsea S.A. & Orbia Advance Corporation S.A.B. de C.V. \\
Coca Cola Femsa S.A.B de C.V. UBL  & \\
Corporacion Inmobiliaria Vesta S.A.B. de C.V. &  \\
El Puerto de Liverpool S.A.B. de C.V. &  \\
Fomento Economico Mexicano S.A.B. de C.V. &  \\
Grupo Carso S.A.B. de C.V. & \\
Genomma Lab Internacional S.A. de C.V. &  \\
Gentera S.A.B. de C.V. & Gruma S.A.B. B \\
Grupo Aeroportuario del Centro Norte S.A.B. de C.V. & \\
Grupo Aeroportuario del Sureste S.A.B. de C.V. B & \\
Grupo Cementos de Chihuahua S.A.B. de C.V. & \\
Grupo Comercial Chedraui S.A. de C.V. & \\
Grupo Mexico S.A.B. de C.V. B & \\
Kimberly Clark de Mexico S.A.B. de C.V. A & \\
La Comer S.A.B. de C.V. UBC & \\
Qualitas Controladora S.A.B. de C.V. &  \\
Regional S.A. de C.V. &  \\
Walmart de Mexico S.A.B. de C.V. & \\
\bottomrule()
\caption{Asignación de los activos acorde a los 2 clústeres }
\end{longtable}

\begin{longtable}[]{@{}
  >{\raggedright\arraybackslash}p{(\columnwidth - 6\tabcolsep) * \real{0.2404}}
  >{\raggedright\arraybackslash}p{(\columnwidth - 6\tabcolsep) * \real{0.2569}}
  >{\raggedright\arraybackslash}p{(\columnwidth - 6\tabcolsep) * \real{0.2408}}
  >{\raggedright\arraybackslash}p{(\columnwidth - 6\tabcolsep) * \real{0.2619}}@{}}
\toprule()
\multicolumn{4}{@{}>{\raggedright\arraybackslash}p{(\columnwidth - 6\tabcolsep) * \real{1.0000} + 6\tabcolsep}@{}}{%
\begin{minipage}[b]{\linewidth}\raggedright
\textbf{Empresas representativas para 4 clústeres}
\end{minipage}} \\
\midrule()
\endhead
\textbf{Bolsa Mexicana de Valores S.A. de C.V.} & \textbf{Genomma Lab Internacional S.A. de C.V.} &
\textbf{Grupo Financiero Banorte O.} & \textbf{Grupo Televisa S.A.B. CPO }\\
Becle S.A. De C.V. & & Grupo Bimbo S.A.B. & Alfa S.A. A.  \\
Industrias Peñoles & & Cemex S.A. C.P.O. & Megacable Holdings S.A.B. de C.V.  \\
La Comer S.A.B de C.V.UBC & & Alsea S.A. & Orbia Advance Corporation S.A.B. de C.V.  \\
& & America Movil S.A.B. de C.V. B & \\
& & Arca Continental S.A.B. de C.V. & \\
& & Banco del Bajio S.A. & \\
& & Coca Cola Femsa S.A.B de C.V. UBL & \\
& & Corporacion Inmobiliaria Vesta S.A.B. de C.V. & \\
& & El Puerto de Liverpool S.A.B. de C.V. & Gruma S.A.B. B \\
& & Fomento Economico Mexicano S.A.B. de C.V.  & \\
& &  Gentera S.A.B. de C.V.  & \\
& & Gruma S.A.B. B & Grupo Comercial Chedraui S.A. de C.V. \\
& & Grupo Aeroportuario del Centro Norte S.A.B. de C.V.  & \\
& & Grupo Aeroportuario del Sureste S.A.B. de C.V. B  & \\
& & Grupo Carso S.A.B. de C.V. & \\
& & Grupo Cementos de Chihuahua S.A.B. de C.V. & \\
& & Grupo Cementos de Chihuahua S.A.B. de C.V. & \\
& & Grupo Financiero Inbursa O. & \\
& & Grupo Mexico S.A.B. de C.V. B. & \\
& & Kimberly Clark de Mexico S.A.B. de C.V. A & \\
& & Qualitas Controladora S.A.B. de C.V. & \\
& & Regional S.A. de C.V. & \\
& & Walmart de Mexico S.A.B. de C.V. & \\
\bottomrule()
\caption{Asignación de los activos acorde a los 4 clústeres }
\end{longtable}

\begin{longtable}[]{@{}
  >{\raggedright\arraybackslash}p{(\columnwidth - 14\tabcolsep) * \real{0.0794}}
  >{\raggedright\arraybackslash}p{(\columnwidth - 14\tabcolsep) * \real{0.1059}}
  >{\raggedright\arraybackslash}p{(\columnwidth - 14\tabcolsep) * \real{0.1312}}
  >{\raggedright\arraybackslash}p{(\columnwidth - 14\tabcolsep) * \real{0.1439}}
  >{\raggedright\arraybackslash}p{(\columnwidth - 14\tabcolsep) * \real{0.1495}}
  >{\raggedright\arraybackslash}p{(\columnwidth - 14\tabcolsep) * \real{0.1289}}
  >{\raggedright\arraybackslash}p{(\columnwidth - 14\tabcolsep) * \real{0.1495}}
  >{\raggedright\arraybackslash}p{(\columnwidth - 14\tabcolsep) * \real{0.1117}}@{}}
\toprule()
\multicolumn{8}{@{}>{\raggedright\arraybackslash}p{(\columnwidth - 14\tabcolsep) * \real{1.0000} + 14\tabcolsep}@{}}{%
\begin{minipage}[b]{\linewidth}\raggedright
\textbf{Empresas representativas para 8 clústeres}
\end{minipage}} \\
\midrule()
\endhead
\textbf{Alfa S.A. A.} & \textbf{America Movil S.A.B. de C.V. B} &
\textbf{Arca Continental S.A.B. de C.V.} & \textbf{Bolsa Mexicana de Valores S.A. de C.V.} & \textbf{Genomma Lab Internacional S.A. de C.V.} &
\textbf{Grupo Carso S.A.B. de C.V. } & \textbf{Grupo Mexico S.A.B. de C.V. B} & \textbf{Grupo Televisa S.A.B. CPO} \\
& Walmart de Mexico S.A.B. de C.V. & Grupo Bimbo S.A.B & Becle S.A. De C.V. &  & Corporacion Inmobiliaria Vesta S.A.B. de C.V. & Cemex S.A. CPO & Megacable Holdings S.A.B. de C.V. \\
& & Banco del Bajio S.A. & Industrias Peñoles &  & Fomento Economico Mexicano S.A.B. de C.V. & Grupo Cementos de Chihuahua S.A.B. de C.V. & Orbia Advance Corporation S.A.B. de C.V. \\
& & Coca Cola Femsa S.A.B de C.V. UBL &  &  & Gentera S.A.B. de C.V. & El Puerto de Liverpool S.A.B. de C.V. & \\
& & Grupo Aeroportuario del Centro Norte S.A.B. de C.V. &  &  & Kimberly Clark de Mexico S.A.B. de C.V. A &  La Comer S.A.B. de C.V. UBC & \\
& & Grupo Aeroportuario del Sureste S.A.B. de C.V. B  & &  & Qualitas Controladora S.A.B. de C.V. & Gruma S.A.B. B & \\
& & Grupo Comercial Chedraui S.A. de C.V.  & &  & Alsea S.A.& & \\
& & Grupo Financiero Banorte O. & & & & & \\
& & Grupo Financiero Inbursa O. & & & &  & \\
& & Regional S.A. de C.V. & & & &  & \\
\bottomrule()
\caption{Asignación de los activos acorde a los 8 clústeres }
\end{longtable}

A primera vista lo que podemos notar de estás tres tablas es que no fueron clústeres con una distribución uniforme, pues en los tres casos hubo uno que notoriamente es conformado por más elementos que el resto y en el caso de 4 y 8 clústeres fue la misma empresa (La Comer S.A.B. de C.V. UBC) la que fue el centroide con más empresas asociadas a él. Cabe destacar que se trata de una empresa comercial y de manera peculiar, podemos observar que su clúster queda integrado por otras empresas comerciales, de servicios y enfocadas a producción.

Esto a diferencia de, por ejemplo, los clústeres de Grupo Financiero Inbursa O. Los cuales están conformado mayormente por empresas financieras. Siendo un caso que comparte Grupo Financiero Banorte O. Es decir, tenemos que las financieras se agruparon con empresas de la misma índole o de rubros similares.

También las tablas 3 y 4, muestras un clúster en el cual se encuentran varias de las empresas más importantes del país en sectores muy distintos. Pues tenemos a Grupo México S.A.B. cd C.V. B (infraestructura, minería y transporte), Grupo Bimbo S.A.B. (alimentos), Grupo Televisa S.A.B. CPO (telecomunicaciones y radiodifusión), Arca Continental S.A.B. de C.V. (bebidas no alcohólicas y alimentos), e Industrias Peñoles (minería, metalurgia, química). Incluso en la tabla 2 están dentro del mismo clúster. 

Hay varias empresas que se mantienen juntas en las tablas, como Grupo Carso S.A.B de C.V. y Bolsa Mexicana de Valores S.A. de C.V. Donde la primera es tal vez la empresa privada más importante del país, especialmente por abarcar los rubros comercial, infraestructura, energía, industrial, construcción. Todos son sectores tales que se encuentra fuertemente vinculado con el gobierno y es muchas veces contratada para cubrir labores estatales. Por eso hace sentido que vaya junto la única empresa completamente nacional del ínidce, como es la Bolsa Mexicana de Valores S.A de C.V.

Por último, cabe resaltar que hay tres clústeres en la tabla 3, los cuales no poseen elementos más que su centroide. Esto podría indicar que la cantidad de centroides debería de ser menor ya que pierde algo de sentido el realizar un clustering para que un clúster tenga 10 elementos y otros 3 solo posean su centroide. En especial siendo que la tabla 3 indica un clustering mejor distribuido. Aunque determinar la cantidad óptima de clústeres se sale de este trabajo

\begin{figure}[htbp]
\centerline{\includegraphics[scale=.5]{Entera/1.png}}
\label{Figura 1}
\caption{Errores semestrales utilizando la capitalización de mercado para determinar los pesos de los activos.}
\end{figure}

\begin{figure}[htbp]
\centerline{\includegraphics[scale=.5]{Entera/2.png}}
\label{Figura 2}
\caption{Errores semestrales utilizando la volatilidad para determinar los pesos de los activos.}
\end{figure}

\begin{figure}[htbp]
\centerline{\includegraphics[scale=.5]{Entera/3.png}}
\label{Figura 3}
\caption{Errores semestrales utilizando el \textit{benchmark equally weighted} para determinar los pesos de los activos.}
\end{figure}

\begin{figure}[htbp]
\centerline{\includegraphics[scale=.5]{Entera/4.png}}
\label{Figura 4}
\caption{Media de los errores cuando se utilizó la capitalización de mercado para determinar los persos de los activos.}
\end{figure}

\begin{figure}[htbp]
\centerline{\includegraphics[scale=.5]{Entera/5.png}}
\label{Figura 5}
\caption{Media de los errores cuando se utilizó la volatilidad para determinar los persos de los activos.}
\end{figure}

La hipótesis de que los \emph{benchmarks} y modelos planteados en trabajos mayormente estadounidenses, no servirían adecuadamente en una economía tan distinta como la mexicana, se vio reflejada en los resultados. Las cinco figuras anteriores nos denotan  esto y además de denotar otros hechos.

Hay dos períodos con mucho más error que el resto, y estos serían el segundo semestre de 2018 y el primero de 2020. En el primero de estos períodos fue el de las elecciones y el cambio de gobierno. Muchos de los cambios de retórica y discursos del gobierno que llegaba al poder provocaron mucha incertidumbre durante dicho período (Bacaria Colom, 2018)\footnote{Bacaria Colom, Jordi (2018). \textit{La nueva presidencia de México: entre la incertidumbre y la esperanza}. Barcelona Center for international affairs.}. Cabe resaltar que si bien durante 2024 también hubo un cambio de gobierno, este solo se presentó como una continuación, mientras que el anterior se vio como un rompimiento total con la alternancia bipartidista de las últimas décadas. Además, en 2024 fue un candidato con mucho más apoyo empresarial y mejor manejo de su discurso que su contraparte en 2018. 

Y durante el primer semestre de 2020, fue el inicio de la pandemia del COVID-19. Etapa también llena de incertidumbre de los mercados e inicio de una recesión global, a causa de la pandemia. Entonces en los períodos donde hubo más error, ocurrieron eventos que nos podrían indicar a qué se deben estos picos en los semestres ya indicados (8 y 11).

Por otro lado, también destacan los períodos, con poco error, es decir, los correspondientes al segundo semestre de 2020, y el segundo de 2021. Durante el primero, hubo una caída del PIB de alrededor 17.8\% (INEGI, 2020)\footnote{INEGI (2020). \textit{Economía y sectores productivos}. [Base de datos]} en el segundo trimestre.  Mientras que en el otro un crecimiento del 1.5\% (INEGI, 2021)\footnote{INEGI (2021). \textit{Economía y sectores productivos}. [Base de datos]} para el término del semestre.  Entonces ambos períodos presentan datos macroeconómicos muy distintos para poder relacionarlos. Lo único que les relaciona es que uno forma parte de las peores caídas tras la pandemia y el otro es el repunte económico cuando esta estaba concluyendo. 

Dados los eventos ocurridos al menos a nivel nacional durante ese período y por lo anteriormente mencionado, no puedo realizar una conjetura apropiada que determine por qué esos períodos fueron notoriamente los de menor error. 

La figura 4 es de interés ya que la culsterización que tuvo menor error, fue la realizada con 2 clústeres. Lo cual no es lo que estamos habituados a intuir, pues por lo general, partimos de la idea que entre más clústeres tengamos, más preciso será nuestro modelo, sin embargo, aquí fue el caso opuesto, a medida que se aumentaron los centroides, aumentó el error. Es también interesante que haya ocurrido utilizando la capitalización de mercado dado que es la manera en que se calcula el IPC.

Tal vez por esto último es que no hayan sido óptimas las otras cantidades de centroides, pues al haber tomado pocos centroides, reducimos el sesgo que pudiese generar el tener una sobrerrepresentación. Podría ser que, bajo el propio criterio de la correlación lineal, y siendo consistente con el uso de la capitalización de mercado, que precisamente esto determinase una cierta robustez del modelo. Podemos decir entonces que la \emph{multilateral insurence} no funciona a partir de una cierta cantidad de clústeres o en un economía como la mexicna.

Una indicación de que tener muchos clústeres, no es garantía de un buen rendimiento del modelo son los resultados en las figuras 4 y 5, pues este nunca es el que presenta menos error, ni se aproxima a este. En particular, cuando se discutió la tabla 10, se habló que había varios centroides sin elementos más que sí mismos, que, junto a los resultados obtenidos por esa métrica, podríamos determinar que la sobrerrepresentación termina por generar un sesgo, en lugar del esperado \emph{multilateral insurence} .

El hecho de no haber agregado una figura para el valor medio presentado por el \textit{benchmark equally weighted} es que los pesos resultantes del \textit{benchmark} de la volatilidad fueron prácticamente idénticos a los otros. Tanto así, que se pueden apreciar las figuras 2 y 3, para notar que se trata prácticamente de la misma. En general no hay alguna diferencia más que en centesimales. Esto significa que los activos seleccionados a través del \textit{clustering} tenían prácticamente la misma volatilidad entre sí.

Esto último podría ser un efecto de la concentración de sectores económicos enteros en pocas empresas, las cuales no son muy diversas. De esta forma, las acciones y activos de las empresas no estarían condicionadas a la competencia y la realidad del mercado nacional, sino que solo se verían afectados por aspectos macroeconómicos que influyen a todas las empresas al mismo tiempo y de maneras similares. 

En general podríamos atribuir el error grande en todas las mediciones a que la manera en que fluctúa el mercado mexicano y sus empresas principales, a que varía mucho más que el de economías desarrolladas, y nuevamente, a que el IPC es medido con tan solo 35 empresas, mientras que los índices de economías desarrolladas son de cien o incluso más como pudieran ser el SP 500 o NASDAQ. Además de que sus índices presentan una mayor diversidad de sectores económicos, así como empresas muy diversificadas y donde una empresa usualmente no concentra la mayoría de una industria. No al menos de la forma en la que ocurre a nivel nacional.

\textbf{Conclusiones}

Le economía mexicana es una muy distinta a la de naciones desarrolladas, y con los resultados que obtuvimos podemos visualizar con mayor claridad cuáles son algunos de estos aspectos, como las prácticas con matices monopólicos o la carencia de una economía diversificada. En cierta medida muchos de estos aspectos pueden verse como áreas de oportunidad.

Ciertamente queda mucho trabajo pendiente por realizar. Es de cierto interés el ver cómo haber implementado propiamente los modelos de Markowitz para haber optimizado los portafolios y reducido el riesgo, y ver qué tanto impacto tendría. Contrastar los datos ya optimizados con estos nos daría una mejor lectura.

Aún así, el que sin haber optimizado, hayamos terminado con errores tan altos es un indicativo claro de que la teoría económica no sirve de manera igual para todos y que la teoría que se utilizó para este trabajo no es tan general como muchos lo quieren hacer ver. 

Ciertamente el encontrar otras formas de medir el error y ver si existen otros acercamientos que nos brinden mejores resultados o se adapten más a nuestra realidad es también una pregunta que queda pendiente. Cabe destacar que dado los resultados atípicos que estaban teniendo, se intentaron realizar otras mediciones y procesos alternativos, pero por falta de conocimiento especializado en el ámbito económico, fueron descartados para mantener un cierto rigor. 

También cabe agregar que, si bien se planteó que se querían probar los modelos durante la recesión de 2008, esto no se realizó al notar que no fue necesario someter al modelo a una situación de ese calibre, puesto que, con el período elegido, se presentaron por si mismos resultados bastante ilustrativos y claros para desestimar que esta clase de modelos y teorías no necesariamente aplica a contextos como el nacional. 


\textbf{Referencias}

\begin{enumerate}
\def\labelenumi{\arabic{enumi}.}
\item
  Gharanchaei, M. K., Panda, P. P. (2024). \emph{Constructing and
  investment fund thourg stock clustering and integer programming}.
  Journal of Advancements in Applied Business Research.
\item
  Markowitz, H. (1952). \emph{Portfolio selection}. JSTOR. (pp 77- 91).
\item
  Wolsey Laurence, A. (1998). Formulations. \emph{Integer Programming}.
  (pp. 1-18). Wiley Interscience.
\item
  Bertsimas, D., Darnell, C., Soucy, R., (1998). \emph{Portfolio construction
  through mixed integer programming}. MIT.
\item
  Kuhn, C. C. N., Calbert, G., Garanovich, I. L., Weir, T., (2023).
  \emph{Integer linear programming supporting portfolio design.}
  Elsevier.
\item
  Anderson, R. M., Bianachi, S. W., Goldberg, L. R. (2011). \emph{Will
  My Risk Parity Strategy Outperform?} Berkeley.
\item
  Baker, M., Bradley, B., Wurgler, J\emph{. Benchmarks as Limits to
  Arbitrage: Understanding the Low-Volatility Anomaly}. Financial
  Analysts Journal.
\item
  Koumou, G. B. (2020). \emph{Diversification and portfolio theory: a
  review}. Swiss Society for Financial Market Research.
\item
  Knutzen, S., Retterholt, J., (2015). \emph{Performance of risk-based
  asset allocation strategies} {[}tesis master, Copenhagen Business
  School{]}.
\item
  Choueifaty, Y., Coignard, Y. (2008). \emph{Toward Maximum
  Diversification}. The Journal of Portfolio Management.
\item
  Reed, T; Muñoz Moreno, R. (febrero 23, 2023). \emph{How the Federal Economic
  Competition Commission fights for workers and consumers in Mexico}.
  World Bank.
\item
  Chow, T., Hsu, J., Kalesnik, V., Little, B. (2011). \emph{A Survey of
  Alternative Equity Index Strategies.} Financial Analysts Journal.
\item
  Carvalho R. L., Lu, X., Moulin, P. (2012). \emph{Demystifying Equity
  Risk-Based Strategies: A Simple Alpha plus Beta Description}. The
  Journal of Portfolio Management.
\item
  Botyarov, M., Miller, E. E., \emph{Partitioning around medoids as a
  systemic approach to generative design solution space reduction.}
  Elsevier.
\item
  S\&P Dow Jones Indices. (febrero 2025\emph{). S\&P/BMV Indices:
  Metodología}.
\item
  Banco de México. (2025). Sistema de Información Económica {[}Base de datos{]}.
\item 
Bacaria Colom, Jordi (2018). \textit{La nueva presidencia de México: entre la incertidumbre y la esperanza}. Barcelona Center for international affairs.
\item 
INEGI (2020). \textit{Economía y sectores productivos}. [Base de datos]
\item 
INEGI (2021). \textit{Economía y sectores productivos}. [Base de datos]
\item
  S\&P Dow Jones Indices. (31 marzo 2025). S\&P/BMV IPC: \emph{Constituents Export}. {[}Base de datos{]}.
\item
  Sánchez Deanny (octubre 6, 2023). Antitrust in Mexico: Perspectives,
  challenges, and opportunities.
\end{enumerate}

\end{document}